\section{The relative categorical family}
\label{sec:category}

We construct an action on bounded complexes of graded projectives and a
family whose stabilizer will eventually be identified with $A_X$.  Fusion
labels realize the possibly nonintegral entries of the Coxeter form by
an ordinary finite-dimensional algebra.  The construction belongs to the
fusion-labelled zigzag framework of Heng and Licata
\cite[Sections~4--5]{HengLicata2024}.  We supply the algebraic details
needed for our relative family, including braid relations as bimodule
chain homotopies.  This strength will allow us to act on right simple
modules as well as on left projective complexes in Section~\ref{sec:caps}.

\subsection{Classical inputs and the relative frame}

Write $A=A_M$.  For $T\subseteq S$, let $W_T$ denote the Coxeter group
defined by the restricted matrix, and call $T$ \emph{spherical} if
$W_T$ is finite.  A positive word means a word in the generators
$\sigma_i$, with no inverses.  Its support is the set of letters that
occur; positive braid relations preserve this set.

\begin{theorem}[Classical Artin and Coxeter inputs]
\label{cat:classical}
The group defined by the Artin presentation on $T\subseteq S$ embeds
in $A$, with image $A_T$.  No sphericity hypothesis is needed here.
For a spherical $T$, the positive Artin monoid $A_T^+$ embeds in
$A_T$, is cancellative, and has greatest common divisors on both
sides.  Every element of $A_T$ is both a left fraction and a right
fraction of positive elements.  Any two elements $a,b\in A_T$ thus
have a common positive extension $pa=qb$, with $p,q\in A_T^+$.

Reduced expressions in any Coxeter system are related by braid moves.
They therefore define a positive lift $\iota(w)$ of a Coxeter element
$w$.  If $T$ is spherical and $w_T$ is its longest element, the left
and right divisors of $\Delta_T=\iota(w_T)$ are precisely
$\{\iota(w):w\in W_T\}$.  In the finite geometric representation,
root signs detect descents, the group acts simply transitively on
chambers, and the translates of a closed fundamental chamber cover
the entire dual space.
\end{theorem}

The embedding statement is
\cite[Chapter~II, Section~4, Theorem~4.13(i)]{VanDerLek1983}; see also
\cite[Section~1, p.~2]{CharneyParis2014}.  For the positive monoid, cancellation,
fractions, and greatest common divisors are
\cite[Propositions~2.3, 4.2, and~5.5; Theorem~5.6]{BrieskornSaito1972}
and \cite[Theorem~4.14]{Deligne1972}; the description of the divisors
of $\Delta_T$ is
\cite[Lemma~5.4 and Propositions~5.5 and~5.7]{BrieskornSaito1972}.
The common-extension assertion follows by writing
$ba^{-1}=q^{-1}p$.  The Coxeter assertions are
\cite[Chapter~IV, \S1; Chapter~V, \S3, no.~3, and \S4,
nos.~4 and~6]{Bourbaki2002}; see also
\cite[Theorems~3.4.2(ii), 4.1.6(i), and Appendix~D]{Davis2008}.
Throughout the proof, positive fractions,
greatest common divisors, and the finite chamber covering are used
only after the relevant subsystem has been shown to be spherical.

Fix $X\subseteq S$.  Add a color $f$ with label $2$ to every color
of $X$ and label $3$ to every color of $S\setminus X$.  Add one
further color $o$, with label $3$ to every color of $S\sqcup\{f\}$,
and put
\[
 \widehat S=S\sqcup\{f,o\}.
\]
Only colors of $S$ will act.  The color $f$ specifies the initial
family, while $o$ is an auxiliary frame in the algebraic construction.
Set $C=2I-b$, where $b$ is symmetric, has zero diagonal, and
\[
 b_{ij}=2\cos(\pi/m_{ij})\quad(m_{ij}<\infty),
 \qquad b_{ij}=2\quad(m_{ij}=\infty).
\]
Thus $b_{if}=0$ for $i\in X$, $b_{if}=1$ for
$i\in S\setminus X$, and $b_{io}=1$ for $i\ne o$.
Two colors are adjacent when $b_{ij}>0$.  The boundary $\partial T$
of a color set consists of its exterior neighbors in $\widehat S$.
All vector inequalities below are coordinatewise.

On root-coefficient vectors, the geometric reflection at $i$ is
\begin{equation}
 s_i v=v-(Cv)_i\varepsilon_i.
 \label{cat:reflection}
\end{equation}
The restriction to a color subsystem is its geometric representation.
The finite-type criterion says that $T$ is spherical exactly when
$C_T$ is positive definite
\cite[Chapter~V, \S4, no.~8, Theorem~2]{Bourbaki2002},
\cite[Theorem~6.12.9]{Davis2008}.

\begin{lemma}[Positive and null Cartan data]
\label{cat:cartan}
For spherical $T$, the inverse $C_T^{-1}$ is nonnegative, and is
strictly positive in each connected block.  Prescribed exterior
coordinates consequently have a unique extension $p$ satisfying
$(Cp)_T=0$.  This extension is nonnegative if the prescribed
coordinates are nonnegative.

If $G$ is nonempty and connected, $v_G>0$, and $C_Gv_G\geq0$, then
either $C_G$ is positive definite or $C_Gv_G=0$.
\end{lemma}

\begin{proof}
For $C_T=2I-b_T>0$, the nonnegative symmetric matrix $b_T$ has
spectral radius less than $2$.  Hence
\[
 C_T^{-1}=\frac12\sum_{n\geq0}(b_T/2)^n,
 \qquad
 p_T=C_T^{-1}b_{T,\widehat S\setminus T}p_{\widehat S\setminus T}.
\]
The first formula gives nonnegativity; a path between two vertices
in one connected block gives a strictly positive entry in some
power, proving strict positivity there.  The second formula gives
the assertions about harmonic extension.

For the final assertion, the following identity holds for every
real vector $u$ on $G$:
\begin{equation}
 u^{\mathsf T}C_Gu
 =\sum_{\substack{i<j\\i,j\in G}}b_{ij}v_iv_j
       \left(\frac{u_i}{v_i}-\frac{u_j}{v_j}\right)^2
   +\sum_{i\in G}\frac{(C_Gv)_i}{v_i}u_i^2.
 \label{cat:scaled-form}
\end{equation}
If one entry of $C_Gv$ is positive, vanishing of the first sum makes
all ratios $u_i/v_i$ equal, and the corresponding term of the second
sum makes that common ratio zero.  The form is then positive
definite.  Otherwise the assumed nonnegative vector $C_Gv$ is zero.
\end{proof}

We call this extension the \emph{harmonic interpolation on $T$}.
Coordinates outside $T$ are not changed by interpolation.

\subsection{Labels and the finite algebra}

We use a finite semisimple spherical $\kk$-linear tensor category
$\cC$ with positive categorical dimensions.  Here is an explicit
choice valid for arbitrary Coxeter labels on a finite set.  For each individual finite
edge of label $m\geq4$, take an independent Temperley--Lieb--Jones
category at circle value $2\cos(\pi/m)$.  Its simple objects
$F_0,\ldots,F_{m-2}$ have dimensions
\[
 \dim F_a=\frac{\sin((a+1)\pi/m)}{\sin(\pi/m)}>0,
 \qquad
 F_1\otimes F_a\cong F_{a-1}\oplus F_{a+1},
\]
where nonexistent endpoint summands are omitted.  These are the
spherical semisimplifications of the planar pairing categories.
For the negligible quotient, finite semisimplicity and spherical
structure, see \cite[Sections~5.1--5.4, especially
Theorems~5.4.4--5.4.5 and the discussion following
Definition~5.4.7]{Chen2014}; the displayed simple list and fusion rule
also appear in \cite[Section~4.1, Equations~(10)--(11)]{HengLicata2024}.  More
explicitly, the Jones--Wenzl idempotents with $a\leq m-2$ have the
positive dimensions displayed above, while the next idempotent has
trace zero and is killed in the semisimplification.  Thus the
resulting category has finitely many simples, not merely finitely
many specified edge objects.

Take the tensor product of these finitely many categories, or finite
dimensional vector spaces if there are no such edges.  For an
oriented edge $(i,j)$ of label $m\geq4$, let $D^1_{ij}$ be $F_1$ in
its factor and the unit in the other factors.  This per-edge $F_1$
choice is a variation of the fusion-label choices in
\cite[Section~4.3, Equation~(15), and Remark~4.4]{HengLicata2024}.
On a label-$3$ edge
take $D^1_{ij}=\mathbf1$, and on an infinite edge take
$D^1_{ij}=\mathbf1\oplus\mathbf1$.  Choose opposite edge objects
as pivotal duals.  Then
\begin{equation}
 \dim_{\cC}D^1_{ij}=b_{ij},\qquad
 \dim_\kk\Hom_{\cC}(L,D^1_{ij}\otimes H)\leq1
 \quad\text{on every finite edge}
 \label{cat:edge-labels}
\end{equation}
for simple $L,H$.  The dimension bound follows from the displayed
fusion rule in a single factor.  No such bound is required on
infinite edges.

Let $\operatorname{Irr}(\cC)$ be a set of simple representatives, and
write
\[
 w_L=\dim_{\cC}L>0,
 \qquad D_0=\sum_{L\in\operatorname{Irr}(\cC)}w_L^2.
\]
We always tensor labels on the right.  No symmetry of the tensor
category is used.  Duality and semisimplicity give the useful identity
\begin{equation}
 \sum_L w_L[D\otimes L:H]=w_H\dim_{\cC}D
 \qquad(H\text{ simple}),
 \label{cat:weighted-labels}
\end{equation}
because the multiplicities on the left are those of $L$ in
$D^*\otimes H$.

For each $i\in\widehat S$ and $D\in\cC$, introduce an additive
projective symbol $P_{i,D}$.  Put
$D^0_{ii}=D^2_{ii}=\mathbf1$, use the edge objects in degree one,
and set all remaining $D^a_{ij}$ to zero.  Define
\begin{equation}
 \Hom^a(P_{i,D},P_{j,E})
       =\Hom_{\cC}(D,D^a_{ij}\otimes E).
 \label{cat:projective-hom}
\end{equation}
Degree-zero maps compose as label maps.  The product of two
degree-one maps is zero unless their colors backtrack; on a
backtrack, compose their label maps and apply the pivotal evaluation
$D^1_{ij}\otimes D^1_{ji}\to\mathbf1$.  Opposite orientations use
the two pivotal evaluations.  All other products of positive lengths
are zero.  This multiplication is associative: compatibility with
degree zero is functoriality, while both bracketings of three
positive lengths vanish.

An \emph{unfolded vertex} is a pair $\alpha=(i,L)$ with $L$ simple;
write $\widetilde T=T\times\operatorname{Irr}(\cC)$.  The finite
basic graded algebra $Z$ is defined by
\[
 e_\alpha Z e_\beta=\Hom(P_\alpha,P_\beta)
\]
with multiplication in path order.  Thus $P_\alpha=Ze_\alpha$ is
a left projective and
$\Hom_Z(P_\alpha,P_\beta)=e_\alpha Z e_\beta$.  Let
$Q_\alpha=e_\alpha Z$ be the right projective and $U_\alpha$ its
simple top.  For $\alpha=(i,L)$, write $t_\alpha$ for the
degree-two element corresponding to $\mathrm{id}_L$.  We call it
the \emph{token}.  The only path lengths are $0,1,2$, and lengths
zero and two each have a one-dimensional space on every diagonal
and vanish off the diagonal.

\begin{lemma}[Trace and adjoint duality]
\label{cat:frobenius}
The functional
\[
 \tau(t_{i,L})=w_L,
 \qquad \tau(Z^0\oplus Z^1)=0
\]
is a symmetric Frobenius trace of degree $-2$.  Opposite arrow
spaces pair perfectly by multiplication followed by $\tau$.  The
radical of $Z$ is $Z^{>0}$, and the socle of every indecomposable
left or right projective is its token line.

For a finite graded bimodule $N$, its adjoint dual satisfies
\begin{equation}
 N^\vee:=\Hom_{Z\text{-left}}(N,Z)
       \cong D_\kk(N)\langle2\rangle,
 \qquad
 (P_\alpha\otimes_\kk Q_\beta)^\vee
       \cong(P_\beta\otimes_\kk Q_\alpha)\langle-2\rangle.
 \label{cat:adjoint-dual}
\end{equation}
This is an involutive contravariant equivalence, fixes the regular
bimodule, and preserves additive radicals.
\end{lemma}

\begin{proof}
Units and tokens pair perfectly.  For arrows, pivotal adjunction
identifies a map $H\to D^1_{ji}\otimes L$ with a map
$D^1_{ij}\otimes H\to L$.  Pairing it against
$L\to D^1_{ij}\otimes H$ is a semisimple trace pairing, nondegenerate
because the trace on each simple matrix block is multiplied by a
nonzero dimension.  The two orders of a backtrack have the same
trace by the pivotal partial-trace identity.  Hence
$\tau(ab)=\tau(ba)$, with no sign from internal degree.

The positive-length ideal is nilpotent and its quotient is the
semisimple diagonal, proving the radical assertion.  A token is
annihilated by the radical.  A nonzero arrow pairs with an opposite
arrow, and a nonzero degree-zero element pairs with a token; neither
can lie in the projective socle.  This proves the socle assertion.

The map $f\mapsto\tau\circ f$ gives the first identification in
\eqref{cat:adjoint-dual}.  A map of degree $k$ gives a functional
of degree $k-2$, which accounts for $\langle2\rangle$.  With the
bimodule convention $(a f b)(n)=f(na)b$, symmetry gives
\[
 \tau(f(na)b)=\tau(bf(na))=\tau(f(bna)),
\]
exactly the vector-dual bimodule action.  The trace identifies
$D_\kk(P_\alpha)=Q_\alpha\langle-2\rangle$ and
$D_\kk(Q_\beta)=P_\beta\langle-2\rangle$.  Reversing the tensor
factors and then applying the shift $2$ proves the second formula.
Shifted vector duality squares to the identity.  Any contravariant
equivalence preserves the additive radical, characterized by
invertibility of $1-gf$ for all opposite maps $g$.
\end{proof}

\subsection{Finite unfoldings}

Let $\widetilde C$ have diagonal entries $2$ and off-diagonal entries
$-\dim_\kk(e_\alpha Z^1e_\beta)$, and set
\begin{equation}
 (\phi v)_{i,L}=w_Lv_i.
 \label{cat:phi}
\end{equation}
The arrow pairing makes $\widetilde C$ symmetric, and
\eqref{cat:weighted-labels} gives
\begin{equation}
 \widetilde C\phi=\phi C.
 \label{cat:intertwining}
\end{equation}
Both statements remain true on every color submatrix.

\begin{lemma}[Finite unfolding]
\label{cat:unfolding}
For spherical $T\subseteq S$, the matrix
$\widetilde C_{\widetilde T}$ is positive definite and simply laced.
Replacing an original reflection $s_i$ by the product of the
commuting reflections in its color gives a homomorphism from $W_T$
to the resulting finite Coxeter group.  Under this homomorphism,
\[
 D_L(\widetilde w)=\widetilde{D_L(w)},
 \qquad
 D_R(\widetilde w)=\widetilde{D_R(w)},
\]
where $D_L$ and $D_R$ are the left and right descent sets.  Every
reduced expression lifts to a reduced expression.  In particular,
the longest element lifts to the longest element.
\end{lemma}

\begin{proof}
The strictly positive vector $v=\phi(C_T^{-1}\mathbf1)$ satisfies
$\widetilde C_{\widetilde T}v=\phi\mathbf1>0$.  Applying
\eqref{cat:scaled-form} to each connected unfolded block proves
positive definiteness on the whole unfolded space, including
directions outside $\phi(\RR^T)$.  An infinite edge cannot occur
inside $T$, since its two-by-two Cartan matrix is singular.
Equation~\eqref{cat:edge-labels} therefore makes every off-diagonal
entry of $\widetilde C_{\widetilde T}$ either $0$ or $-1$.  The
finite Coxeter criterion applies to this simply laced matrix.

There are no arrows between distinct vertices of one color.  The
corresponding reflections commute, and their product acts on the
folded space as $s_i$, by \eqref{cat:intertwining}.  Choose $p$ with
$C_Tp>0$.  Then $\phi p$ is strictly inside a full unfolded
fundamental chamber.  For an arbitrary word representing $w$, let $g$
be its lifted product; then $g\phi p=\phi(wp)$.  A root has coefficients either
all nonnegative or all nonpositive in the simple-root basis, so
pairing it with a chamber-interior point detects its sign.  The
row at $(i,L)$ of $\widetilde C\phi(wp)$ is $w_L(Cwp)_i$.
These row signs are the signs of the pairings with
$g^{-1}\epsilon_{(i,L)}$ and $w^{-1}\epsilon_i$, respectively.
Consequently the unfolded left descents are exactly all vertices
over the original left descents.  Applying the same reasoning to
inverse words proves the right-descent assertion.

A word representing the identity fixes $\phi p$.  The stabilizer
of a chamber-interior point in a finite Coxeter group is trivial,
so its lift is the identity.  Thus the lifted product depends only
on $w$.  If an original reduced expression is extended by an ascent,
every vertex of the added color is an ascent of the lifted prefix.
Acting at one of those vertices does not change the signs at the
others because they are mutually orthogonal.  The whole block is
therefore a sequence of ascents.  Induction proves reducedness of
the lift.  The longest original element has every simple descent;
its lift has every unfolded simple descent, so is the longest
unfolded element.
\end{proof}

\subsection{Twists as bimodule homotopies}

The evaluation and coevaluation complexes below adapt the zigzag
bimodule twists of \cite[Proposition~2.4 and Theorem~2.5]{KhovanovSeidel2002}
and \cite[Section~4.6, Proposition~7]{HuerfanoKhovanov2001}; compare also
\cite[Definition~5.2 and Proposition~5.3]{HengLicata2024}.
We keep the full bimodule homotopies, because later tensor products
will involve nonprojective right simples.

All complexes are bounded cochain complexes of finite internally
graded modules.  The differential has cochain degree $1$ and
internal degree $0$.  The shift $\langle a\rangle$ raises internal
placements by $a$, while $[1]$ lowers cochain placements by $1$.
Only cochain degrees contribute signs to tensor products and Hom
complexes.

Set $M_\alpha=P_\alpha\otimes_\kk Q_\alpha$, with multiplication
$\mu_\alpha:M_\alpha\to Z$.  Choose homogeneous dual bases
$p_a\in P_\alpha$, $q_a\in Q_\alpha$ for the trace pairing, so
$\tau(q_ap_b)=\delta_{ab}$.  The element
$\sum_a p_a\otimes q_a$ is central and has internal degree $2$:
centrality is the invariance of the perfect pairing under the
two $Z$ actions.  It defines a bimodule map
$\delta_\alpha:Z\to M_\alpha\langle-2\rangle$.  Put
\begin{equation}
 B_\alpha=[M_\alpha\xrightarrow{\mu_\alpha}Z],
 \qquad
 D_\alpha=[Z\xrightarrow{\delta_\alpha}
                      M_\alpha\langle-2\rangle],
 \label{cat:elementary-twists}
\end{equation}
with $Z$ in cochain degree zero.  Thus the projective terms lie
in cochain degrees $-1$ and $1$, respectively.

\begin{proposition}[The categorical action]
\label{cat:action}
There are bimodule chain homotopy equivalences
$B_\alpha D_\alpha\simeq Z\simeq D_\alpha B_\alpha$.
For a color $i\in S$, define
\[
 B_i=\prod_{\alpha\in\widetilde{\{i\}}}B_\alpha,
 \qquad
 D_i=\prod_{\alpha\in\widetilde{\{i\}}}D_\alpha.
\]
These products are independent of their order up to bimodule
homotopy, are mutually inverse, and satisfy all the defining Artin
relations.  They define an action of $A$ by invertible graded
bimodule homotopy classes, denoted $B_x$ for $x\in A$.
In particular,
\begin{equation}
 F_i(E):=B_i\otimes_ZE
  =\Cone\left(
       \bigoplus_{\alpha\in\widetilde{\{i\}}}
       P_\alpha\otimes_\kk e_\alpha E\longrightarrow E\right).
 \label{cat:color-evaluation}
\end{equation}
Every $B_x$ preserves bounded projective complexes and preserves
the cohomology of their bigraded Hom complexes.  Its adjoint dual
$B_x^\vee$ represents its inverse.  All the action identities are
actual bimodule chain homotopy equivalences and remain valid after
tensoring with any finite module complex, whether or not projective.

For $H\in\cC$, right tensoring of labels defines an additive functor
\[
 T_H(P_{i,D})=P_{i,D\otimes H}
\]
on graded projectives.  It preserves path lengths, extends to
bounded homotopy categories, and commutes naturally up to homotopy
with $B_x\otimes_Z-$ for every $x\in A$.
\end{proposition}

\begin{proof}
We first verify the inverse explicitly.  The product
$B_\alpha D_\alpha$ has $M_\alpha$ in degree $-1$,
\[
 Z\oplus P_\alpha\otimes_\kk(e_\alpha Ze_\alpha)
                          \otimes_\kk Q_\alpha\langle-2\rangle
\]
in degree zero, and $M_\alpha\langle-2\rangle$ in degree $1$.
Write $\theta_\alpha=t_\alpha/w_L$ when $\alpha=(i,L)$, so that
$\tau(\theta_\alpha)=1$.  Split the middle factor into its unit
and $\theta_\alpha$ lines.  The incoming differential onto the
token summand is an isomorphism, up to the tensor sign: its
coefficient is the dual-basis identity
\[
 \sum_a\tau(qp_a)q_a=q.
\]
The token has degree $2$, so this summand is unshifted
$M_\alpha$.  Cancel that isomorphism.  The remaining unit summand
maps identically onto $M_\alpha\langle-2\rangle$ by multiplication
with $e_\alpha$; cancel it as well.  Only the regular bimodule $Z$
remains.  In the opposite product $D_\alpha B_\alpha$, the
dual-basis identity gives the same incoming cancellation, followed
by the unit cancellation.  These cancellations take place in
bimodule complexes, not just after evaluating on projectives.

We next prove the elementary braid relation using the twist-conjugation
argument of \cite[Lemma~2.11 and Propositions~2.12--2.13]{SeidelThomas2001}.
For a bounded perfect
left complex $P$, let its evaluation twist be the bimodule cone
\[
 T_P=\Cone\bigl(P\otimes_\kk\Hom_Z(P,Z)\longrightarrow Z\bigr).
\]
It is unchanged by an internal or cochain shift of $P$.  A homotopy
equivalence of $P$ identifies its evaluation map, and hence its
cone, up to homotopy.  If $B$ is a product of the elementary
invertible complexes, each of its terms is projective on both
sides.  Tensor--Hom adjunction identifies its tensor adjoint with
$B^\vee=\Hom_Z(B,Z)$; the evaluation and coevaluation identities
hold at chain level.  Since $B$ is invertible, this adjoint is its
inverse in the bimodule homotopy category.  Lemma~\ref{cat:frobenius}
keeps its terms in the same class.  Tensor--Hom and evaluation now
give
\begin{equation}
 B T_P B^{-1}\simeq T_{BP}.
 \label{cat:twist-conjugation}
\end{equation}

If distinct vertices $\alpha,\beta$ have no arrows between them,
$M_\alpha M_\beta=M_\beta M_\alpha=0$, so both
$B_\alpha B_\beta$ and $B_\beta B_\alpha$ are the simultaneous
evaluation cone with projective term $M_\alpha\oplus M_\beta$.
If there is exactly one arrow in each direction, the complexes
\[
 B_\alpha P_\beta
  =[P_\alpha\langle1\rangle\longrightarrow P_\beta],
 \qquad
 D_\beta P_\alpha
  =[P_\alpha\longrightarrow P_\beta\langle-1\rangle]
\]
have degrees $(-1,0)$ and $(0,1)$.  Both maps are nonzero by the
perfect backtrack pairing, and their common arrow space is
one-dimensional.  Rescaling one term gives
\[
 B_\alpha P_\beta\simeq
       (D_\beta P_\alpha)\langle1\rangle[1].
\]
Using $T_{P_\gamma}=B_\gamma$ in
\eqref{cat:twist-conjugation} therefore gives
$B_\alpha B_\beta D_\alpha\simeq D_\beta B_\alpha B_\beta$.
Multiplication on the left by $B_\beta$ and on the right by
$B_\alpha$ proves
\[
 B_\alpha B_\beta B_\alpha
       \simeq B_\beta B_\alpha B_\beta.
\]

Vertices within one color are orthogonal.  Their twists commute,
their inverse products are the $D_i$, and all mixed terms in the
expansion vanish.  This proves the simultaneous cone
\eqref{cat:color-evaluation}.  For $m_{ij}<\infty$, the two
alternating color words of length $m_{ij}$ are reduced expressions
for the longest element of the spherical rank-two subsystem.
Lemma~\ref{cat:unfolding} lifts both to reduced expressions for
the same longest unfolded element.  Matsumoto's Theorem relates
them by commutations and three-term braid moves, which we have
just proved as bimodule homotopies.  For $m_{ij}=\infty$ there is
no relation to impose.  This proves the action for arbitrary labels.

It remains to verify the natural commutation with labels.  On
\eqref{cat:projective-hom}, $T_H$ tensors the label map on the
right by $\mathrm{id}_H$.  Evaluation of an edge and its dual
always occurs to the left of the label, so composition and path
length are preserved.  On $P_{j,D}$, the evaluation term for
color $i$ has the canonical semisimple description
\[
 \bigoplus_{a,L}P_{i,L}\langle a\rangle\otimes_\kk
             \Hom_{\cC}(L,D^a_{ij}\otimes D)
 \cong\bigoplus_a P_{i,D^a_{ij}\otimes D}\langle a\rangle.
\]
Its evaluation is the tautological label map.  Tensoring by $H$
replaces $D$ by $D\otimes H$ in this description, naturally in
$D$ and in every projective morphism.  Thus $T_H$ commutes with
the positive evaluation cone.  Composing this natural
identification with the inverse equivalence on both sides proves
the inverse commutation without requiring $T_H$ to be invertible.

Invertible functors preserve all morphism spaces, including both
shifts, and hence preserve bigraded Hom cohomology between
projective complexes.  Finally, tensoring an explicit chain
homotopy with any module complex preserves its defining equations
with the usual cochain signs.  No flatness assertion about the
additional module complex is needed.
\end{proof}

\subsection{The transported family and folded counts}

For $x\in A$ and a simple label $L$, define
\begin{equation}
 E_L(x)=B_xP_{f,L}.
 \label{cat:family-definition}
\end{equation}
The forthcoming layer construction is summarized here to specify
the data used throughout the proof.  A minimal projective complex
has no path-length-zero entries in its differential.  A generator in
cochain placement $p$ and internal placement $a$ has \emph{level}
$p+a$.  A length-$j$ differential entry changes level by $1-j$.
The pure layer $Y^d(E)$ takes the generators of level $d$ and the
length-one differential.  Lemma~\ref{lay:structure} proves the
existence and uniqueness of minimal models and the functorial layer
and truncation properties used below.  Dimensions of pure layers
always count their projective generators, not the dimensions of
the underlying $Z$-modules.

\begin{proposition}[Relative-family properties]
\label{cat:family}
There are nonnegative folded vectors $m^d(x)$ characterized by
\begin{equation}
 \phi m^d(x)=\sum_Lw_L\dim Y^d(E_L(x)).
 \label{cat:folded-counts}
\end{equation}
They have finite level support.  Each individual coordinate ranges
over a locally finite subset of $\RR_{\geq0}$.  Weighted counts of
tops and socles of the layers are likewise proportional within a
color: a nonzero top or socle in that color occurs at every
unfolded vertex of the color in some member of the family.
The same proportionality holds for families of subobjects obtained
functorially by twists and layer truncations.

If only colors of $T\subseteq S$ act, all generator spaces with
their cochain and internal placements outside $\widetilde T$ stay
unchanged up to isomorphism.  In particular, each $E_L(x)$ has
exactly its original $f$-generator in placements $(0,0)$, and
\[
 m_f^0(x)=1,\qquad m_f^d(x)=0\quad(d\ne0).
\]
Every $E_L(x)$ is indecomposable in the graded homotopy category.
The color support of its minimal complex is connected and contains
$f$.
\end{proposition}

\begin{proof}
Proposition~\ref{cat:action} gives
$E_L(x)\simeq T_L(E_{\mathbf1}(x))$.  Tensoring labels preserves
the positive path length of a minimal differential, and therefore
preserves minimality.  If a multiplicity object $D$ occurs at a
fixed color and pair of placements in the unit-label complex, its
contribution in the $L$-member is $D\otimes L$.  For each target
simple $H$, Equation~\eqref{cat:weighted-labels} says that the
weighted multiplicity is $w_H\dim_{\cC}D$.  Summing over placements
of one level proves \eqref{cat:folded-counts}.

Lemma~\ref{lay:labels} proves that label tensoring commutes with
layers, truncations, and the top and socle tests.  Its proof uses
the exactness of tensoring in the semisimple category, together
with \eqref{cat:projective-hom} and
\eqref{cat:weighted-labels}.  It gives the stated proportionality
for functorially constructed subobjects and for tops and socles.
Positivity of all label weights then proves the assertion about
occurrence at every vertex of a color.

There are finitely many labels and finitely many terms in each
complex.  To see local finiteness, fix one target label $H$.
For a bound $R$, Equation~\eqref{cat:folded-counts} expresses
$w_Hm_i^d\leq w_HR$ as a nonnegative integer combination of the
finitely many fixed positive numbers $w_L$.  Each coefficient is
bounded separately.  There are thus only finitely many possible
values below $R$.

For an unfolded vertex $\beta$ outside color $i$,
$U_\beta e_\alpha=0$ for every $\alpha$ in that color.  The
explicit complexes \eqref{cat:elementary-twists} consequently give
$U_\beta B_i=U_\beta=U_\beta D_i$.  Tensoring a minimal left
complex with $U_\beta$ kills its radical differential, leaving
exactly its generator spaces at $\beta$ in their two placements.
This proves exterior invariance, and in particular the statement
about the $f$-generator and its folded counts.

The initial projective $P_{f,L}$ is indecomposable, and an
equivalence preserves indecomposability, proving the first final
assertion.  A disconnected color support would split its minimal
complex into the summands on different components: a differential
uses only arrows, which join adjacent colors, and tokens, which
stay at one color.  Every nonzero minimal summand is
noncontractible, since reducing a putative contracting homotopy
equation modulo the radical would give an identity equal to zero.
Such a splitting contradicts indecomposability.  Exterior
invariance ensures that $f$ belongs to the support.
\end{proof}
